\clearpage\thispagestyle{empty}
\pdfbookmark[0]{Quellen zu Mondsterne (Februar 1823)}{nicolai1823-sources}
\begin{center}\textbf{Quellen zu Mondsterne (Februar 1823)}\end{center}
\noindent Diese KI-Übertragung ins heutige Deutsch folgt Nicolais \textit{Mondsterne} in \href{https://gdz.sub.uni-goettingen.de/id/PPN236007467}{Gauß’ \textit{Werken}, Band~VI} (Göttingen, 1874), S.~443--445. Der Aufsatz erschien zuerst als \textit{Berechnung der Meridiandifferenz zweier Orte, aus correspondirenden Mondsculminationen} in den von H.~C. Schumacher herausgegebenen \textit{Astronomischen Nachrichten}, Band~II, Nr.~26 (Februar 1823), Sp.~\href{https://api.digitale-sammlungen.de/iiif/image/v2/bsb10049427_00021/full/1800,/0/default.jpg}{17--18}, \href{https://api.digitale-sammlungen.de/iiif/image/v2/bsb10049427_00022/full/1800,/0/default.jpg}{19--20}, \href{https://api.digitale-sammlungen.de/iiif/image/v2/bsb10049427_00025/full/1800,/0/default.jpg}{21--22} und \href{https://api.digitale-sammlungen.de/iiif/image/v2/bsb10049427_00026/full/1800,/0/default.jpg}{23--24}. Die Links führen zum Exemplar der Bayerischen Staatsbibliothek. Der Band erschien 1824 in Altona.
\par\smallskip\noindent Die Werkausgabe enthält die Überschrift \textit{Mondsterne} und die eingeklammerte Adressierung an Schumacher. Ihre Reihe aus 22 Strichen kennzeichnet die Auslassung von Nicolais Rechenbeispielen aus Sp.~20--24 des Erstdrucks. Dazu gehören Beobachtungen vom 3., 4. und 5. März und ein Vergleich zwischen Dorpat und Mannheim. Die Übertragung folgt der Notation der Werkausgabe, einschließlich der Reihenfolge der Faktoren im Ausdruck, der die stündliche Bewegung mit der Meridiandifferenz verknüpft, des Wurzelzeichens ohne Überstrich, des Buchstabens $s$ im Absatz nach der Formel für den wahrscheinlichen Fehler und des hochgestellten aufrechten s in $0^{\mathrm{s}}1$. Im Erstdruck steht dagegen in Sp.~24 $0^{\prime\prime}{,}1$. KI-Übertragung und KI-Satz.
\clearpage
